%! Justifying a Claim About the Difference of Two Means Based on a CI — Unit 7, Lesson 7.7 — Group Activity (Answer Key)
\documentclass[10pt]{article}
\usepackage{apstats-article}
\usepackage{apstats-key}   % pulls in -boxes; do NOT also load -boxes
\newcommand{\TallMath}[1]{$\displaystyle #1\rule[-1.4em]{0pt}{3.2em}$}

\begin{document}

\pageheader{Unit 7, Lesson 7.7}{Group Activity --- Answer Key}
\namepartnerperiod

\vspace{0.1in}

\begin{scenariobox}[Tier R --- Interpreting a Given Interval (Template Task)]{sky}
A two-sample $t$-interval for $\mu_{\text{city}} - \mu_{\text{suburban}}$ daily commute time
(minutes) is reported as a 90\% CI: $(4.2,\; 12.6)$ minutes.

\medskip
\textbf{1.} Fill in the interpretation sentence using the template below.

``We are \ans{90}\% confident that the true difference in \ans{population mean daily commute time}
(\ans{city} $-$ \ans{suburban}) commute time is between \ans{4.2} and
\ans{12.6} minutes.''

\medskip
\textbf{2.} Is 0 inside or outside the interval? Circle:
\quad \textbf{Inside} \quad / \quad \ans{Outside}

\medskip
\textbf{3.} Does the interval provide evidence that city residents have a longer mean commute
time than suburban residents? Circle one and explain.

\quad \ans{Yes} \qquad \textbf{No}

\ansline{The entire interval is above 0, so 0 is not plausible; city residents have a higher mean commute time.}
\end{scenariobox}

\begin{scenariobox}[Tier A --- Calculate, Interpret, and Justify]{sky}
A nutrition researcher compares mean daily fiber intake (grams) for adults in two age groups.
Group 1 (ages 18--34): $\bar x_1 = 18.4$ g, $s_1 = 5.2$ g, $n_1 = 45$.
Group 2 (ages 35--54): $\bar x_2 = 21.9$ g, $s_2 = 6.1$ g, $n_2 = 50$.
Both conditions are met. $df = 88$ (given). $t^* = 1.988$ for 95\% confidence.

\textbf{a.} Compute $SE = \sqrt{\dfrac{s_1^2}{n_1} + \dfrac{s_2^2}{n_2}}$.

\[SE = \sqrt{\dfrac{{\color{keyred}\mathbf{5.2}}^2}{{\color{keyred}\mathbf{45}}} + \dfrac{{\color{keyred}\mathbf{6.1}}^2}{{\color{keyred}\mathbf{50}}}} \approx \ans{1.085}\text{ g}\]

\textbf{b.} Construct the 95\% CI for $\mu_1 - \mu_2$.

\[({\color{keyred}\mathbf{18.4}} - {\color{keyred}\mathbf{21.9}}) \pm {\color{keyred}\mathbf{1.988}} \cdot {\color{keyred}\mathbf{1.085}} \approx (\ans{-5.66},\; \ans{-1.34})\text{ g}\]

\textbf{c.} Interpret the interval in context.
\ansline{We are 95\% confident that the true difference in population mean daily fiber intake (ages 18--34 minus ages 35--54) is between $-5.66$ and $-1.34$ grams.}

\textbf{d.} A dietitian claims adults aged 18--34 consume less fiber than those aged 35--54.
Does the CI support this claim? Explain using the interval.

\ansline{Yes. The entire interval is negative (below 0), indicating $\mu_1 < \mu_2$; this supports the claim that younger adults consume less fiber on average.}
\end{scenariobox}

\begin{scenariobox}[Tier E --- Compare Widths and Design a Study]{sky}
Researchers study sleep duration (hours) for college students at two campuses.
Both groups have $s_1 = s_2 = 1.5$ hr. They plan to use 95\% confidence with $t^* \approx 1.97$.

\textbf{a.} Compute the margin of error for each sample-size plan below.

\renewcommand{\arraystretch}{1.8}
\begin{tabular}{@{} l c c c @{}}
  \textbf{Plan} & $n_1 = n_2$ & $SE = \sqrt{\frac{s_1^2}{n_1}+\frac{s_2^2}{n_2}}$ & ME $= t^* \cdot SE$ \\[2pt] \hline
  X & 25  & \ans{$\approx 0.424$ hr} & \ans{$\approx 0.836$ hr} \\
  Y & 100 & \ans{$\approx 0.212$ hr} & \ans{$\approx 0.418$ hr} \\
\end{tabular}

\medskip
\textbf{b.} By what factor did the margin of error decrease from Plan X to Plan Y?
\ansline{The ME decreased by a factor of $\approx 2$ (i.e., $1/\sqrt{4}$), because sample sizes quadrupled.}

\textbf{c.} A researcher wants the margin of error to be no more than 0.30 hours.
Using the same $s$ values and $t^* \approx 1.97$, what minimum equal sample size $n$
is needed in each group? Set up the inequality and solve.

\ansline{$1.97 \cdot \sqrt{1.5^2/n + 1.5^2/n} \le 0.30$; $\sqrt{4.5/n} \le 0.152$; $4.5/n \le 0.0231$; $n \ge 195$.}

\textbf{d.} If the 95\% CI for $\mu_1 - \mu_2$ under Plan Y is $(-0.1,\; 0.5)$ hr,
does this provide evidence that the two campuses differ in mean sleep? Justify.

\ansline{No. 0 is inside the interval $(-0.1, 0.5)$, so a difference of 0 is plausible; insufficient evidence of a difference.}
\end{scenariobox}

\end{document}
